% Generated by officeParser.
\documentclass[aspectratio=169,12pt,xcolor=table]{beamer}
\usepackage{iftex}
\ifPDFTeX
  \usepackage[T1]{fontenc}
  \usepackage[utf8]{inputenc}
  \usepackage{lmodern}
\else
  \usepackage{fontspec}
\fi
\usepackage{array}
\usepackage[normalem]{ulem}
\hypersetup{colorlinks=true,
  linkcolor=blue,
  urlcolor=blue,
  citecolor=blue,
  pdfcreationdate={D:20251125163603Z},
  pdfmoddate={D:20251207101557Z},
  pdfinfo={TestString={Hello from custom props},TestNumber={42},TestBool={true}}}
\setbeamertemplate{frametitle continuation}{}
\definecolor{hex000000}{HTML}{000000}
\definecolor{hex0070C0}{HTML}{0070C0}
\definecolor{hex17365D}{HTML}{17365D}
\definecolor{hex365F91}{HTML}{365F91}
\definecolor{hex4BACC6}{HTML}{4BACC6}
\definecolor{hex4F81BD}{HTML}{4F81BD}
\definecolor{hex7F7F7F}{HTML}{7F7F7F}
\definecolor{hex808080}{HTML}{808080}
\definecolor{hex92D050}{HTML}{92D050}
\definecolor{hex9BBB59}{HTML}{9BBB59}
\definecolor{hexD3DFEE}{HTML}{D3DFEE}
\definecolor{hexFF0000}{HTML}{FF0000}
\definecolor{hexFFFF00}{HTML}{FFFF00}
\definecolor{hexFFFFFF}{HTML}{FFFFFF}

\begin{document}

\begin{frame}[allowframebreaks]
\frametitle{\protect\textcolor{hex17365D}{Demonstration of DOCX support in~calibre}}
\phantomsection\label{demonstration-of-docx-support-in-calibre}
This document~demonstrates~the ability of the~calibre~DOCX Input plugin to convert the various typographic features in a Microsoft Word (2007 and newer) document. Convert this document to a modern~ebook~format, such as AZW3 for Kindles or EPUB for other~ebook~readers, to see it in action.

There is support for images, tables, lists, footnotes, endnotes,~links,~dropcaps~and~various~types~of text and paragraph level formatting.

To see the DOCX conversion in action, simply add this file to~calibre~using the~\textbf{“Add Books”~}button and then click “\textbf{Convert”.~}~Set the output format in the top right corner of the conversion dialog to EPUB or AZW3 and click~\textbf{“OK”}.
\end{frame}
\note{Slide Note 1.

1}

\begin{frame}[allowframebreaks]
\frametitle{\protect\textcolor{hex365F91}{Text Formatting}}
\phantomsection\label{text-formatting}
\textcolor{hex4F81BD}{\textbf{{\fontsize{13pt}{15.6pt}\selectfont Inline formatting}}}

Here, we~demonstrate~various types~of inline text formatting and the use of embedded fonts.

Here~is~some~\textbf{bold,~}\textit{italic,~}\textit{\textbf{bold-italic,~}}\uline{underlined~}and~\sout{struck~out~}~text. Then, we have a super\textsuperscript{script}~and a sub\textsubscript{script}. Now we see some~\textcolor{hexFF0000}{red},~\textcolor{hex92D050}{green}~and~\textcolor{hex0070C0}{blue}~text. Some text with a~{\setlength{\fboxsep}{1pt}\colorbox{hexFFFF00}{\strut yellow} \colorbox{hexFFFF00}{\strut highlight.}} Some text in a~box. Some text~in~\textcolor{hexFFFFFF}{{\setlength{\fboxsep}{1pt}\colorbox{hex000000}{\strut inverse} \colorbox{hex000000}{\strut video}}}.

A paragraph with styled text:~\textcolor{hex808080}{\textit{subtle~emphasis~~}}followed~by~\textbf{strong text~}and~\textcolor{hex4F81BD}{\textit{\textbf{intense emphasis}}}.~This paragraph uses document wide styles for styling rather than inline text properties as~demonstrated~in the~previous~paragraph~—~calibre~can handle both with equal ease.

\textcolor{hex4F81BD}{\textbf{{\fontsize{13pt}{15.6pt}\selectfont Fun with fonts}}}

This document has embedded the Ubuntu font family. The body text is in the Ubuntu typeface, here is~\texttt{some text in the Ubuntu Mono typeface, notice how every letter has the same width, even~i~and m}. Every embedded font will automatically be embedded in the output~ebook~during conversion.~

\textcolor{hex4F81BD}{\textbf{{\fontsize{13pt}{15.6pt}\selectfont Paragraph level formatting}}}

{\raggedleft You can do crazy things with~paragraphs, if~the urge strikes you. For~instance~this paragraph is~right aligned~and has a right border. It has also been given a light gray background.\par}

For~the lovers~of poetry amongst you, paragraphs with hanging indents, like this often come in handy. You can use hanging indents to ensure that a line of poetry~retains~its individual identity as a line even when the screen~is~ too~narrow to display it as a single line. Not only does this paragraph have a hanging indent,~it~is also has an extra top margin, setting it apart from the preceding paragraph.
\end{frame}
\note{Slide Note 2. This is a \textbf{bold} text, but no \textit{italic}. Let's strike \sout{that} out.

2}

\begin{frame}[allowframebreaks]
\frametitle{\protect\textcolor{hex365F91}{Tables}}
\phantomsection\label{tables}
\begin{tabular}{|*{2}{>{\raggedright\arraybackslash}p{\dimexpr(\linewidth-4\tabcolsep-3\arrayrulewidth)/2\relax}|}}
\hline
\cellcolor{hex9BBB59}\textcolor{hexFFFFFF}{\textbf{ITEM~}} & \cellcolor{hex9BBB59}\textcolor{hexFFFFFF}{\textbf{NEEDED~}} \\
\hline
Books~ & 1~ \\
\hline
Pens~ & 3~ \\
\hline
Pencils~ & 2~ \\
\hline
Highlighter~ & 2 colors~ \\
\hline
Scissors~ & 1 pair~ \\
\hline
\end{tabular}

Tables in Word can vary from the extremely simple to the extremely complex.~calibre~tries to do~its~best when converting tables. While you may run into trouble with the occasional table,~the vast majority of~common cases should be converted very well, as~demonstrated~in this section.~Note that for~optimum~results, when creating tables in Word, you should set their widths using percentages, rather than absolute units.~~To the left of this paragraph is a floating two column table with a nice green border and header row.~

Now~let’s~look at a fancier table—one with alternating row colors and partial borders. This table is stretched out to take 100\% of the available width.

{\small\setlength{\tabcolsep}{4pt}
\begin{tabular}{|*{6}{>{\raggedright\arraybackslash}p{\dimexpr(\linewidth-12\tabcolsep-7\arrayrulewidth)/6\relax}|}}
\hline
\cellcolor{hexFFFFFF}{\fontsize{11pt}{13.2pt}\selectfont City or Town~} & \cellcolor{hexFFFFFF}{\fontsize{11pt}{13.2pt}\selectfont Point A~} & \cellcolor{hexFFFFFF}{\fontsize{11pt}{13.2pt}\selectfont Point B~} & \cellcolor{hexFFFFFF}{\fontsize{11pt}{13.2pt}\selectfont Point C~} & \cellcolor{hexFFFFFF}{\fontsize{11pt}{13.2pt}\selectfont Point D~} & \cellcolor{hexFFFFFF}{\fontsize{11pt}{13.2pt}\selectfont Point E~} \\
\hline
\cellcolor{hexFFFFFF}Point A~ & \cellcolor{hexD3DFEE}—~ & \cellcolor{hexD3DFEE}~ & \cellcolor{hexD3DFEE}~ & \cellcolor{hexD3DFEE}~ & \cellcolor{hexD3DFEE}~ \\
\hline
\cellcolor{hexFFFFFF}Point B~ & 87~ & —~ & ~ & ~ & ~ \\
\hline
\cellcolor{hexFFFFFF}Point C~ & \cellcolor{hexD3DFEE}64~ & \cellcolor{hexD3DFEE}56~ & \cellcolor{hexD3DFEE}—~ & \cellcolor{hexD3DFEE}~ & \cellcolor{hexD3DFEE}~ \\
\hline
\cellcolor{hexFFFFFF}Point D~ & 37~ & 32~ & 91~ & —~ & ~ \\
\hline
\cellcolor{hexFFFFFF}Point E~ & \cellcolor{hexD3DFEE}93~ & \cellcolor{hexD3DFEE}35~ & \cellcolor{hexD3DFEE}54~ & \cellcolor{hexD3DFEE}43~ & \cellcolor{hexD3DFEE}—~ \\
\hline
\end{tabular}
}
\end{frame}
\note{And then about tables

3}

\begin{frame}[allowframebreaks]
\relax
Next, we see a table with special formatting in various locations. Notice how the formatting for the header row and sub header rows is preserved.

\begin{tabular}{|*{4}{>{\raggedright\arraybackslash}p{\dimexpr(\linewidth-8\tabcolsep-5\arrayrulewidth)/4\relax}|}}
\hline
\cellcolor{hex4BACC6}\textcolor{hexFFFFFF}{\textbf{College~}} & \cellcolor{hex4BACC6}\textcolor{hexFFFFFF}{\textbf{New students~}} & \cellcolor{hex4BACC6}\textcolor{hexFFFFFF}{\textbf{Graduating students~}} & \cellcolor{hex4BACC6}\textcolor{hexFFFFFF}{\textbf{Change~}} \\
\hline
~ & \textcolor{hex808080}{\textit{Undergraduate}}~ & ~ & ~ \\
\hline
Cedar University~ & {\fontsize{11pt}{13.2pt}\selectfont 110~} & {\fontsize{11pt}{13.2pt}\selectfont 103~} & {\fontsize{11pt}{13.2pt}\selectfont +7~} \\
\hline
Oak Institute~ & {\fontsize{11pt}{13.2pt}\selectfont 202~} & {\fontsize{11pt}{13.2pt}\selectfont 210~} & {\fontsize{11pt}{13.2pt}\selectfont -8~} \\
\hline
~ & \textcolor{hex808080}{\textit{Graduate}}~ & ~ & ~ \\
\hline
Cedar University~ & {\fontsize{11pt}{13.2pt}\selectfont 24~} & {\fontsize{11pt}{13.2pt}\selectfont 20~} & {\fontsize{11pt}{13.2pt}\selectfont +4~} \\
\hline
Elm College~ & {\fontsize{11pt}{13.2pt}\selectfont 43~} & {\fontsize{11pt}{13.2pt}\selectfont 53~} & {\fontsize{11pt}{13.2pt}\selectfont -10~} \\
\hline
\cellcolor{hexFFFFFF}Total~ & \cellcolor{hexFFFFFF}{\fontsize{11pt}{13.2pt}\selectfont 998~} & \cellcolor{hexFFFFFF}{\fontsize{11pt}{13.2pt}\selectfont 908~} & \cellcolor{hexFFFFFF}{\fontsize{11pt}{13.2pt}\selectfont 90~} \\
\hline
\end{tabular}

\textcolor{hex808080}{\textit{{\fontsize{10pt}{12pt}\selectfont Source:}}}{\fontsize{10pt}{12pt}\selectfont ~Fictitious data, for illustration purposes only~}

Next, we have something a little more complex, a nested table,~i.e.~a table inside another table. Additionally, the inner table has some of its cells merged.~The table is displayed horizontally centered.~

\begin{tabular}{|*{2}{>{\raggedright\arraybackslash}p{\dimexpr(\linewidth-4\tabcolsep-3\arrayrulewidth)/2\relax}|}}
\hline
One~

Three~ & Two~ \\
\hline
 & Four~ \\
\hline
\end{tabular}

\begin{tabular}{|*{2}{>{\raggedright\arraybackslash}p{\dimexpr(\linewidth-4\tabcolsep-3\arrayrulewidth)/2\relax}|}}
\hline
~ & To the left is a table inside a table, with some cells merged.~ \\
\hline
\end{tabular}
\end{frame}
\note{Now calendars

\begin{enumerate}
\item Sun
\item Mon
\end{enumerate}

7. Sat

5}

\begin{frame}[allowframebreaks]
\relax
We~end~with a~fancy~calendar, note how~much of the original formatting is preserved.~Note that this table will only display correctly on~relatively wide~screens. In general, very wide tables or tables whose cells have fixed width requirements~don’t~fare well in~ebooks.

{\scriptsize\setlength{\tabcolsep}{2pt}
\begin{tabular}{|*{14}{>{\raggedright\arraybackslash}p{\dimexpr(\linewidth-28\tabcolsep-15\arrayrulewidth)/14\relax}|}}
\hline
\textcolor{hex365F91}{December 2007~} &  &  &  &  &  &  &  &  &  &  &  &  &  \\
\hline
\textcolor{hex365F91}{Sun~} & ~ & \textcolor{hex7F7F7F}{Mon~} & ~ & \textcolor{hex7F7F7F}{Tue~} & ~ & \textcolor{hex7F7F7F}{Wed~} & ~ & \textcolor{hex7F7F7F}{Thu~} & ~ & \textcolor{hex7F7F7F}{Fri~} & ~ & \textcolor{hex7F7F7F}{Sat~} &  \\
\hline
~ & ~ & ~ & ~ & ~ & ~ & ~ & ~ & ~ & ~ & ~ & ~ & \textcolor{hex7F7F7F}{1~} &  \\
\hline
~ & ~ & ~ & ~ & ~ & ~ & ~ & ~ & ~ & ~ & ~ & ~ & ~ &  \\
\hline
\textcolor{hex365F91}{2~} & ~ & \textcolor{hex7F7F7F}{3~} & ~ & \textcolor{hex7F7F7F}{4~} & ~ & \textcolor{hex7F7F7F}{5~} & ~ & \textcolor{hex7F7F7F}{6~} & ~ & \textcolor{hex7F7F7F}{7~} & ~ & \textcolor{hex7F7F7F}{8~} &  \\
\hline
~ & ~ & ~ & ~ & ~ & ~ & ~ & ~ & ~ & ~ & ~ & ~ & ~ &  \\
\hline
\textcolor{hex365F91}{9~} & ~ & \textcolor{hex7F7F7F}{10~} & ~ & \textcolor{hex7F7F7F}{11~} & ~ & \textcolor{hex7F7F7F}{12~} & ~ & \textcolor{hex7F7F7F}{13~} & ~ & \textcolor{hex7F7F7F}{14~} & ~ & \textcolor{hex7F7F7F}{15~} &  \\
\hline
~ & ~ & ~ & ~ & ~ & ~ & ~ & ~ & ~ & ~ & ~ & ~ & ~ &  \\
\hline
\textcolor{hex365F91}{16~} & ~ & \textcolor{hex7F7F7F}{17~} & ~ & \textcolor{hex7F7F7F}{18~} & ~ & \textcolor{hex7F7F7F}{19~} & ~ & \textcolor{hex7F7F7F}{20~} & ~ & \textcolor{hex7F7F7F}{21~} & ~ & \textcolor{hex7F7F7F}{22~} &  \\
\hline
~ & ~ & ~ & ~ & ~ & ~ & ~ & ~ & ~ & ~ & ~ & ~ & ~ &  \\
\hline
\textcolor{hex365F91}{23~} & ~ & \textcolor{hex7F7F7F}{24~} & ~ & \textcolor{hex7F7F7F}{25~} & ~ & \textcolor{hex7F7F7F}{26~} & ~ & \textcolor{hex7F7F7F}{27~} & ~ & \textcolor{hex7F7F7F}{28~} & ~ & \textcolor{hex7F7F7F}{29~} &  \\
\hline
~ & ~ & ~ & ~ & ~ & ~ & ~ & ~ & ~ & ~ & ~ & ~ & ~ &  \\
\hline
\textcolor{hex365F91}{30~} & ~ & \textcolor{hex7F7F7F}{31~} & ~ & ~ & ~ & ~ & ~ & ~ & ~ & ~ & ~ & ~ &  \\
\hline
\end{tabular}
}
\end{frame}

\begin{frame}[allowframebreaks]
\frametitle{\protect\textcolor{hex365F91}{Structural Elements}}
\phantomsection\label{structural-elements}
Miscellaneous structural elements you can add to your document, like footnotes, endnotes,~dropcaps~and the like.~

\textcolor{hex4F81BD}{\textbf{{\fontsize{13pt}{15.6pt}\selectfont Footnotes \& Endnotes}}}

Footnotes\textsuperscript{1}~and endnotes\textsuperscript{i}~are automatically recognized and both are converted to endnotes, with backlinks for maximum ease of use in~ebook~devices.

\textcolor{hex4F81BD}{\textbf{{\fontsize{13pt}{15.6pt}\selectfont Dropcaps}}}

{\fontsize{59pt}{70.8pt}\selectfont D}

rop~caps are used to emphasize the leading paragraph at the start of a section. In~Word~it is possible to specify how many lines of text a drop-cap should use.

\textcolor{hex4F81BD}{\textbf{{\fontsize{13pt}{15.6pt}\selectfont Links}}}

Two kinds of links are possible, those that refer to an external~website~and those that refer to~locations~inside the document itself. Both are supported by~calibre. For example, here is a link pointing to the~\href{http://calibre-ebook.com/download}{\uline{calibre download page}}.~Then we have a link that points back to the section on~\uline{paragraph level formatting}~in this document.
\end{frame}

\begin{frame}[allowframebreaks]
\frametitle{\protect\textcolor{hex365F91}{Images}}
\phantomsection\label{images}
Centered images like this are useful for large pictures that should be a focus of attention.~

% alt: 165 Inspirational Quotes To Keep You Motivated In Life
\includegraphics[width={\ifdim 810pt>\linewidth\linewidth\else 810pt\fi},height={\ifdim 810pt>0.75\textheight0.75\textheight\else 810pt\fi},keepaspectratio]{images/image1.png}

There is no analogous technology in~ebooks, so~theconversion~will usually end up placingthe image either centered or floating close to the point in the text where itwas~\textcolor{hex808080}{\textit{inserted}},~notnecessarily~where it appears on the page in Word.
\end{frame}

\begin{frame}[allowframebreaks]
\relax
{\centering \textcolor{hex365F91}{\textbf{{\fontsize{14pt}{16.8pt}\selectfont Lists}}}\par}

\textcolor{hex365F91}{All types of lists are supported by the conversion,~with the exception of~lists that use fancy~bullets,~these get converted to regular bullets.}

\textcolor{hex4F81BD}{\textbf{{\fontsize{13pt}{15.6pt}\selectfont Bulleted List}}}

\begin{itemize}
\item \textcolor{hex000000}{One}
\item \textcolor{hex000000}{Two}
\end{itemize}

\textcolor{hex4F81BD}{\textbf{{\fontsize{13pt}{15.6pt}\selectfont Numbered List}}}

\begin{enumerate}
\item One, with~a very long~line to~demonstrate~that the hanging indent for the list is working correctly
\item Two
\end{enumerate}

\textcolor{hex4F81BD}{\textbf{{\fontsize{13pt}{15.6pt}\selectfont Multi-level Lists}}}

\begin{enumerate}
\item One
\begin{enumerate}
\item Two
\begin{enumerate}
\item Three
\item Four with~a very long~line to~demonstrate~that the hanging indent for the list is working correctly.
\item Five
\end{enumerate}
\end{enumerate}
\item Six
\end{enumerate}

A~Multi-level~list with bullets:

\begin{itemize}
\item One
\begin{itemize}
\item Two
\begin{itemize}
\item This bullet uses an image as the bullet item
\item Four
\end{itemize}
\end{itemize}
\item Five
\end{itemize}

\textcolor{hex4F81BD}{\textbf{{\fontsize{13pt}{15.6pt}\selectfont Continued Lists}}}

\begin{enumerate}
\item One
\item Two
\end{enumerate}

An interruption in our regularly scheduled listing, for this essential and~very relevant~public service announcement.

\begin{enumerate}
\setcounter{enumi}{2}
\item We now resume our normal programming
\item Four
\end{enumerate}
\end{frame}

\begin{frame}[allowframebreaks]
\frametitle{Charts}
\phantomsection\label{charts}
\begin{tabular}{|*{4}{>{\raggedright\arraybackslash}p{\dimexpr(\linewidth-8\tabcolsep-5\arrayrulewidth)/4\relax}|}}
\hline
 & Series 1 & Series 2 & Series 3 \\
\hline
Category 1 & 4.3 & 2.4 & 2 \\
\hline
Category 2 & 2.5 & 4.4000000000000004 & 2 \\
\hline
Category 3 & 3.5 & 1.8 & 3 \\
\hline
Category 4 & 4.5 & 2.8 & 5 \\
\hline
\end{tabular}
\end{frame}

\end{document}
